\documentclass[10pt,a4paper]{amsart}
\usepackage{amsmath,amssymb,amsthm,graphicx}
\usepackage{verbatim}
\newenvironment{reference}{\comment}{\endcomment}
%\newenvironment{reference}{}{}
\newenvironment{slogan}{\comment}{\endcomment}
\newenvironment{history}{\comment}{\endcomment}

% For commutative diagrams we use Xy-pic
\usepackage[all]{xy}

% We use 2cell for 2-commutative diagrams.
\xyoption{2cell}
% Theorem environments.
%
\theoremstyle{plain}
\newtheorem{theorem}[subsection]{Theorem}
\newtheorem{proposition}[subsection]{Proposition}
\newtheorem{lemma}[subsection]{Lemma}

\theoremstyle{definition}
\newtheorem{definition}[subsection]{Definition}
\newtheorem{example}[subsection]{Example}
\newtheorem{exercise}[subsection]{Exercise}
\newtheorem{situation}[subsection]{Situation}

\theoremstyle{remark}
\newtheorem{remark}[subsection]{Remark}
\newtheorem{remarks}[subsection]{Remarks}

\numberwithin{equation}{subsection}

% Macros
%
\def\lim{\mathop{\mathrm{lim}}\nolimits}
\def\colim{\mathop{\mathrm{colim}}\nolimits}
\def\Spec{\mathop{\mathrm{Spec}}}
\def\Hom{\mathop{\mathrm{Hom}}\nolimits}
\def\Ext{\mathop{\mathrm{Ext}}\nolimits}
\def\SheafHom{\mathop{\mathcal{H}\!\mathit{om}}\nolimits}
\def\SheafExt{\mathop{\mathcal{E}\!\mathit{xt}}\nolimits}
\def\Sch{\mathit{Sch}}
\def\Mor{\mathop{\mathrm{Mor}}\nolimits}
\def\Ob{\mathop{\mathrm{Ob}}\nolimits}
\def\Sh{\mathop{\mathit{Sh}}\nolimits}
\def\NL{\mathop{N\!L}\nolimits}
\def\CH{\mathop{\mathrm{CH}}\nolimits}
\def\proetale{{pro\text{-}\acute{e}tale}}
\def\etale{{\acute{e}tale}}
\def\QCoh{\mathit{QCoh}}
\def\Ker{\mathop{\mathrm{Ker}}}
\def\Im{\mathop{\mathrm{Im}}}
\def\Coker{\mathop{\mathrm{Coker}}}
\def\Coim{\mathop{\mathrm{Coim}}}

% Boxtimes
%
\DeclareMathSymbol{\boxtimes}{\mathbin}{AMSa}{"02}

%
% Macros for moduli stacks/spaces
%
\def\QCohstack{\mathcal{QC}\!\mathit{oh}}
\def\Cohstack{\mathcal{C}\!\mathit{oh}}
\def\Spacesstack{\mathcal{S}\!\mathit{paces}}
\def\Quotfunctor{\mathrm{Quot}}
\def\Hilbfunctor{\mathrm{Hilb}}
\def\Curvesstack{\mathcal{C}\!\mathit{urves}}
\def\Polarizedstack{\mathcal{P}\!\mathit{olarized}}
\def\Complexesstack{\mathcal{C}\!\mathit{omplexes}}
% \Pic is the operator that assigns to X its picard group, usage \Pic(X)
% \Picardstack_{X/B} denotes the Picard stack of X over B
% \Picardfunctor_{X/B} denotes the Picard functor of X over B
\def\Pic{\mathop{\mathrm{Pic}}\nolimits}
\def\Picardstack{\mathcal{P}\!\mathit{ic}}
\def\Picardfunctor{\mathrm{Pic}}
\def\Deformationcategory{\mathcal{D}\!\mathit{ef}}
\title{Distinct monoid characters into a field are linearly independent (Artin's theorem)}
\author{The Stacks Project Authors}
\date{stacks-project:fields.tex:local-snapshot-20261012; GNU FDL 1.2 or later}
\begin{document}
\maketitle
\noindent\textit{Editorial note.} All delivered mathematical text is the author's own, copied from the pinned Stacks Project chapter file fields.tex; no statement or proof step was rewritten, reordered or removed.\par\medskip
\noindent\textit{Editorial note.} The counted proof is the author's complete argument: The coefficient-vanishing branch is compressed by the author into 'Since lambda\_i not = 0 we conclude that chi\_n(h) = chi\_i(h) for all i by induction' (source lines 1709-1711): the displayed relation holds for all g, so induction on the n-1 distinct characters forces each coefficient lambda\_i (chi\_n(h) - chi\_i(h)) to vanish.\par\medskip
\noindent\textit{Editorial note.} The label lemma-independence-characters names the counted statement, and the proof environment immediately following it in the source is delivered verbatim as its proof.\par\medskip
\noindent\textit{Editorial note.} The counted statement is carried by the author's own lemma environment in the chapter source.\par\medskip
\noindent\textit{Allowed standard prerequisites (editorial).}\begin{itemize}
\item Monoid homomorphisms G -> L (L regarded as a monoid under multiplication) send the neutral element of G to 1
\item Same chapter use: the Artin-Schreier and power-sum arguments both invoke this lemma (not part of the counted unit)
\end{itemize}\medskip
\begin{lemma}
\label{lemma-independence-characters}
Let $L$ be a field. Let $G$ be a monoid, for example a group. Let
$\chi_1, \ldots, \chi_n : G \to L$ be pairwise distinct
homomorphisms of monoids where $L$ is regarded as a monoid
by multiplication. Then $\chi_1, \ldots, \chi_n$
are $L$-linearly independent: if $\lambda_1, \ldots, \lambda_n \in L$
not all zero, then $\sum \lambda_i\chi_i(g) \not = 0$
for some $g \in G$.
\end{lemma}

\begin{proof}
If $n = 1$ this is true because $\chi_1(e) = 1$ if $e \in G$ is the
neutral (identity) element. We prove the result by induction for $n > 1$.
Suppose that $\lambda_1, \ldots, \lambda_n \in L$ not all zero.
If $\lambda_i = 0$ for some, then we win by induction on $n$.
Since we want to show that $\sum \lambda_i\chi_i(g) \not = 0$
for some $g \in G$ we may after dividing by $-\lambda_n$
assume that $\lambda_n = -1$. Then the only way we get in trouble
is if
$$
\chi_n(g) = \sum\nolimits_{i = 1, \ldots, n - 1} \lambda_i\chi_i(g)
$$
for all $g \in G$. Fix $h \in G$. Then we would also get
\begin{align*}
\chi_n(h)\chi_n(g) & = \chi_n(hg) \\
& = \sum\nolimits_{i = 1, \ldots, n - 1} \lambda_i\chi_i(hg) \\
& = \sum\nolimits_{i = 1, \ldots, n - 1} \lambda_i\chi_i(h) \chi_i(g)
\end{align*}
Multiplying the previous relation by $\chi_n(h)$ and subtracting we obtain
$$
0 = \sum\nolimits_{i = 1, \ldots, n - 1}
\lambda_i (\chi_n(h) - \chi_i(h)) \chi_i(g)
$$
for all $g \in G$. Since $\lambda_i \not = 0$ we conclude that
$\chi_n(h) = \chi_i(h)$ for all $i$ by induction.
The choice of $h$ above was arbitrary, so we conclude
that $\chi_i = \chi_n$ for $i \leq n - 1$ which contradicts
the assumption that our characters $\chi_i$ are pairwise distinct.
\end{proof}

\end{document}
